\begin{abstract}
Learning to coordinate in long-horizon cooperative tasks is difficult when
agents share a reward that obscures their individual contributions. Feudal
hierarchies structure these tasks through goals assigned by a manager to
workers, but shared feedback does not distinguish useful goal assignments
from progress produced by teammates. We introduce a counterfactual-goal
framework for credit assignment in Feudal multiagent reinforcement learning
with simple spatial waypoint goals. A manager assigns each worker a waypoint
for a fixed horizon, while workers learn control policies using an intrinsic
reward based on reducing their distance to the assigned waypoint. Both levels
are trained with proximal policy optimization. To assign credit to the
manager's decisions, a learned advantage model evaluates alternative
waypoints for one worker while holding the other workers' waypoints fixed.
The resulting counterfactual baseline adjusts the team advantage to provide
a separate learning signal for each worker's goal assignment. This connects
credit assignment across agents and hierarchical levels while retaining an
explicit, interpretable goal definition. We evaluate the approach on a
tightly coupled continuous-control multiagent task involving cooperative
object pushing, where successful manipulation requires coordinated physical
interactions among agents. Our results show that our algorithm can scale to larger agent counts and handle sparse rewards better than the baselines.
% TODO: Replace the placeholder with validated findings for this method.
\textit{[Results pending: specify the comparison methods and principal
empirical finding, including a quantitative result.]}
\end{abstract}
